% !TEX root = 78-circularity-audit.tex
% ASSERT_CONVENTION: natural_units=natural, metric_signature=mostly_minus, fourier_convention=NA, coupling_convention=NA, renormalization_scheme=NA, gauge_choice=NA
%==============================================================================
% Phase 78 (v18.0 Cartan/MacDowell-Mansouri): Circularity Audit -- forced vs posited
% THE FINAL phase of v18.0. This file = the complete Phase-78 derivation.
% Plan 78-01 (Sections 1-3): THE DECISIVE forced-vs-posited computation:
%   (1) the bare so(3,1)-invariant quad-in-curvature 4-form space dim=2 anchor
%       (Euler eps R^R + Pontryagin R^R), exact over Q;
%   (2) Tr(X o Y)|frame: the bare-trace (4,0) Fubini-Study FOIL vs the soldered (1,3)
%       eta -- the trace form supplies the symmetric frame metric (=> the Pontryagin
%       source, and up to orientation a volume sqrt|eta| eps);
%   (3) THE DECISIVE trace-form-invariant subspace dimension + eps-reachability (both
%       routes) + det_3-fixed-vs-free normalization -> the decisive triple.
%   The input ban (no posited MM/EH action) is declared and AST/source-guarded.
% Plan 78-02 (Sections 4-7, this completion): the verdict ladder applied to the triple
%   (Def. ladder + Prop. verdict => fp-imported-action, with the named shortfalls); the
%   Lambda=0 corollary (eps F^F -> Gauss-Bonnet topological at the measured vacuum); the
%   GST/Singh/Castro imported-action contrast class (where each imports its action); and
%   the milestone-closing v18.0 summary (74 foundation / 75 SURVIVES coframe / 76
%   reinterpreted Berry / 77 NEGATIVE matter-Einstein / 78 forced-vs-posited) + the
%   combined milestone verdict. The verdict is reported at TRUE STRENGTH (neither
%   inflated into a win nor deflated out of conservatism) and is BLOCKING-human-ratified
%   (78-02 Task 3; not self-ratified in this file).
%==============================================================================
\documentclass[11pt]{article}
\usepackage{amsmath,amssymb,amsthm}
\usepackage[margin=1in]{geometry}
\newcommand{\R}{\mathbb{R}}
\newcommand{\Q}{\mathbb{Q}}
\newcommand{\OO}{\mathbb{O}}
\newcommand{\half}{\tfrac12}
\newcommand{\eu}{e}
\newcommand{\om}{\omega}
\newcommand{\Tr}{\operatorname{Tr}}
\newcommand{\so}{\mathfrak{so}}
\newcommand{\hh}{\mathfrak{h}}
\DeclareMathOperator{\Pf}{Pf}
\theoremstyle{definition}\newtheorem{remark}{Remark}
\theoremstyle{definition}\newtheorem{definition}{Definition}
\theoremstyle{plain}\newtheorem{proposition}{Proposition}
\theoremstyle{plain}\newtheorem{lemma}{Lemma}

\title{Phase 78: The MacDowell--Mansouri $\varepsilon$-contraction is \emph{not}
forced by the $\hh_3(\OO)$ trace form / cubic norm --- the forced-vs-posited
circularity audit and the milestone-closing verdict (exact over $\Q$)}
\author{v18.0 --- Gravity as the Curvature of the Peirce-Frame Cartan/MM Connection on $\hh_3(\OO)$}
\date{2026-06-02 (78-01 decisive count); 2026-06-03 (78-02 verdict)}

\begin{document}
\maketitle

\begin{abstract}
We compute, exact over $\Q$, the decisive integer that settles whether the
MacDowell--Mansouri (MM) $\varepsilon$-contraction --- the $\mathrm{Spin}(9,1)\to
\mathrm{SO}(3,1)$ symmetry breaking plus the $\varepsilon$-tensor that turns the
Cartan curvature term $F\wedge F$ into Einstein--Hilbert $+\,\Lambda$ --- is
\emph{forced} by the intrinsic $\hh_3(\OO)$ trace form $\langle X,Y\rangle=\Tr(X\circ
Y)$ and cubic norm $\det\nolimits_3$, or appears \emph{only} because an MM/EH action
was posited by hand ($\texttt{fp-imported-action}$). The audited object is the
\emph{linear-in-Riemann} Einstein--Hilbert term
$\varepsilon_{abcd}R^{ab}\wedge \eu^c\wedge \eu^d$ (the cross-term of the
$\varepsilon F\wedge F$ square), \emph{not} a quadratic-in-$F$ Pontryagin/Maxwell
stress (the overturned Phase-76 tautology). We find: (i) the bare
$\so(3,1)$-invariant quadratic-in-curvature $4$-form space is exactly
$2$-dimensional (Euler $+$ Pontryagin); (ii) the trace form supplies the symmetric
frame metric $\eta$, hence the Pontryagin contraction is reachable; (iii)
$\varepsilon$ is reachable \emph{only} as the metric volume form $\sqrt{|\det\eta|}\,
\varepsilon$, whose orientation is a discrete choice not fixed by the symmetric
$\Tr/\det\nolimits_3$ data, and the cubic norm $\det\nolimits_3$ \emph{vanishes
identically} on the soldered Lorentz $4$-block so it supplies no orientation and no
normalization. The decisive triple is therefore
$\bigl(\dim=1\text{ (}\eta\text{-tensorial, Pontryagin only)}\,/\,2\text{ (admitting
the volume form, }\varepsilon\text{ non-unique})\,;\ \varepsilon\text{-in-span }=\,$
YES only via the volume form$\,;\ \text{normalization }=$ free/imported$\bigr)$ ---
the input to the $\texttt{fp-imported-action}$ verdict, reported at true strength.
Applying the verdict ladder (Plan~78-02), the milestone verdict is
$\texttt{fp-imported-action}$ on all three counts (the $1$-dimensional intrinsic
subspace is generated by Pontryagin not $\varepsilon$; $\varepsilon$'s orientation is
not intrinsically fixed; the normalization is free) --- an honest partial in the same
equivalence class as the GST / Singh / Castro posited-action programs, reinforced by the
$\Lambda=0$ corollary (at the Phase-77-measured vacuum even a posited
$\varepsilon$-action is the topological Gauss--Bonnet term). Together with Phase~77's
NEGATIVE dynamical verdict ($G[g]$ not Einstein-form intrinsically), this closes
milestone v18.0: the $\hh_3(\OO)$ Peirce-frame Cartan/MM connection gives a curved,
matter-sourced, position-dependent Lorentzian slice carrying a forced
$\mathrm{SO}(3,1)$ coframe, but not Einstein gravity without a posited action. The
final verdict is rendered deterministically by the driver and ratified at the blocking
human checkpoint.
\end{abstract}

\section*{Conventions (locked; CONVENTIONS.md \S0,\S1,\S6,\S11; state.json \texttt{convention\_lock})}
Everything decisive is \textbf{exact over $\Q$} (\texttt{sympy.Matrix.rank}/
\texttt{nullspace} over \texttt{QQ}); never \texttt{numpy.linalg}, never float on the
decisive dimension/rank/identity ($\texttt{fp-float-decisive}$). The det SSOT is
\texttt{ring\_lemma\_verification.py} $\det\nolimits_3$ --- the \emph{unique}
$F_4$-invariant cubic norm (cross-term $2\,\mathrm{Re}((x_2x_1)x_3)$, $x_2$ before
$x_1$; $324/324$ inner-derivation annihilation, Phase~64.1/65) ---
\texttt{octonion\_algebra.py} is banned. The trace form is
$c(X,Y)=\Tr(X\circ Y)=\Tr(\mathrm{jordan}(X,Y))$, $\mathrm{jordan}=\half(AB+BA)$,
$F_4$-invariant. Signature mostly-minus $(+,-,-,-)$, frame metric
$\eta=\mathrm{diag}(+1,-1,-1,-1)$. Gravity is the Lorentz block $R[\om]$ of
$F=dA+A\wedge A$, $A=\om\oplus(1/\ell)\eu$, with (Wise, gr-qc/0611154)
\begin{equation}\label{eq:wise}
  F=\Bigl(R[\om]-\tfrac{\Lambda}{3}\,\eu\wedge \eu\Bigr)+d_\om \eu .
\end{equation}
The residual structure group is $\so(3,1)[6]$, \emph{forced} by $(E_{11},u)$
(Phase~75 SURVIVES; residual $21=\so(3,1)[6]\oplus\so(6)[15]$, the $\so(6)$ a
trivial-on-spacetime ideal). The invariant count is at the \emph{broken}
$\so(3,1)$ level: the $\mathrm{Spin}(9,1)\to\mathrm{SO}(3,1)$ breaking is exactly the
``by hand'' step under audit. At the Phase-77-measured vacuum $\Lambda=0$ the
$\varepsilon F\wedge F$ action is the pure Gauss--Bonnet/Euler topological term;
we do not reintroduce $\Lambda<0$.

\paragraph{The audited object (anti-tautology; \texttt{fp-wrong-object}).}
Substituting $F^{ab}=R^{ab}-\tfrac{\Lambda}{3}\eu^a\wedge \eu^b$ into
$\varepsilon_{abcd}F^{ab}\wedge F^{cd}$ and expanding,
\begin{equation}\label{eq:ehexpand}
  \varepsilon_{abcd}F^{ab}\wedge F^{cd}
  =\underbrace{\varepsilon_{abcd}R^{ab}\wedge R^{cd}}_{\text{Gauss--Bonnet (topological)}}
   -\tfrac{2\Lambda}{3}\underbrace{\varepsilon_{abcd}R^{ab}\wedge \eu^c\wedge \eu^d}_{\text{Einstein--Hilbert (LINEAR in }R)}
   +\tfrac{\Lambda^2}{9}\,\varepsilon_{abcd}\eu^a\wedge \eu^b\wedge \eu^c\wedge \eu^d .
\end{equation}
The Einstein--Hilbert term is \emph{linear} in the Riemann curvature $R[\om]$, wedged
with the tetrad; it is \emph{not} the quadratic $R\wedge R$ Pontryagin piece. The
$\varepsilon$ whose forced-vs-posited status we audit is the one producing this
linear-in-$R$ term. (At $\Lambda=0$ only $R\wedge R$ survives, manifestly topological
--- a second, independent reason the audit lands on $\texttt{fp-imported-action}$.)

\medskip\noindent\emph{Driver:} \texttt{code/cartan\_phaseC\_contraction.py}
(exact over $\Q$; reuses the warm engines \texttt{ring\_lemma\_verification},
\texttt{cartan\_phaseA\_coframe}, \texttt{cartan\_phaseB\_curvature},
\texttt{orbit\_dimension\_gate}; no \texttt{octonion\_algebra}, no
\texttt{numpy.linalg} on the decisive path). All numerals below are reproduced by
that driver.

%==============================================================================
\section{Task 1 --- the certified intrinsic toolkit and the input ban}
%==============================================================================

\subsection{The det SSOT re-pass}
We re-pass the cheap algebraic identities of the cubic norm (the $F_4$-uniqueness
certificate being the $324/324$ inner-derivation annihilation of Phase~64.1/65,
cited): $\det\nolimits_3(I)=1$; $\det\nolimits_3(\mathrm{diag}(a,b,c))=abc$
(symbolic over $\Q$); and the polarization headline
\begin{equation}
  d(X,X,X)=6\,\det\nolimits_3(X)\qquad(\text{exact over }\Q),
\end{equation}
with $d(X,Y,Z)=N(X{+}Y{+}Z)-N(X{+}Y)-N(X{+}Z)-N(Y{+}Z)+N(X)+N(Y)+N(Z)$ the
symmetric trilinear polarization of $N=\det\nolimits_3$. The driver confirms
$\det\nolimits_3(\mathrm{diag}(2,3,5))=30$ and $\texttt{octonion\_algebra}\notin
\texttt{sys.modules}$ on the decisive path.

\subsection{The two frame trace-form metrics: the $(4,0)$ FOIL and the soldered $(1,3)$ $\eta$}
There are two distinct ``trace-form'' metrics on the $4$-frame, and keeping them apart
is essential.

\begin{remark}[the Euclidean Fubini--Study foil, transparency only]
The \emph{bare} Jordan trace Gram on the raw frame directions
$\{\beta,\gamma,p,q\}=$ engine indices $[1,2,3,10]$ is
\begin{equation}
  \bigl[\Tr(\mathrm{jordan}(\eu_i,\eu_j))\bigr]=\mathrm{diag}(1,1,2,2),
  \qquad\text{signature }(4,0).
\end{equation}
This is the compact $\mathrm{OP}^2=F_4/\mathrm{Spin}(9)$ Fubini--Study metric (the
real part of the QGT). It is reported as the diagnostic FOIL and is \emph{never} the
spacetime metric (Phase~75).
\end{remark}

\begin{lemma}[the soldered frame metric is the $\det_2$ Lorentzian $\eta$]\label{lem:eta}
The physical frame metric is the soldering bilinear $B(\delta,\delta')=(\delta\circ
\delta')|_{V_0}$ projected onto $V_0\cong\R^{3,1}$, equivalently the $\det_2$ form
$\det\nolimits_2=x_0^2-x_1^2-x_2^2-x_3^2$ in the $52$-kkt Minkowski coordinates
$x^0=(\beta+\gamma)/2,\ x^1=p,\ x^2=q,\ x^3=(\beta-\gamma)/2$. Its raw Gram
$G_{\rm DET2,RAW}$ has eigenvalues $\{\half,-\half,-1,-1\}$, signature $(1,3)$, and in
the orthonormal frame equals $\eta=\mathrm{diag}(+1,-1,-1,-1)$. The change of basis
$T$ ($v_{\rm on}=Tv_{\rm raw}$) satisfies $G_{\rm DET2,RAW}=T^{\!\top}\eta\,T$ (exact
over $\Q$).
\end{lemma}

It is this soldered $\eta$ --- the symmetric frame metric supplied by the (projected)
trace/Jordan structure --- that gives (a) the Pontryagin contraction
$\eta^{ac}\eta^{bd}$ and (b), \emph{up to orientation}, a volume form
$\sqrt{|\det\eta|}\,\varepsilon$. Everything in Tasks~2--3 is in this orthonormal
frame; the invariant-tensor dimension is conjugation-invariant, so the decisive
integer is frame-independent.

\subsection{The hard input ban}
The forced determination may use \emph{only} $\Tr(X\circ Y)$, $\det\nolimits_3$, the
Peirce decomposition under $E_{11}$, the $\mathbb{C}_u/\pi_u$ reduction, the
$(E_{11},u)$-forced $\so(3,1)$ structure, and standard differential geometry. The
driver's \texttt{input\_ban\_guard()} (mirroring
\texttt{cartan\_phaseB\_einstein.ast\_guard\_T}/\texttt{source\_guard}) asserts that
\emph{no} load-bearing posited MM/EH action, no $\int\varepsilon F\wedge F$ as an
assumed action, no ``$-1/2$''/``$16\pi G$'' coupling, no SUSY/MM symbol, no
\texttt{octonion\_algebra}, and no \texttt{numpy.linalg} appears on the decisive path
(AST for identifiers; source-string with docstring/comment/string-literal stripping
for coefficient phrases; runtime for \texttt{sys.modules}). The guard passes on the
decisive functions and still catches a genuine load-bearing violation (verified).

%==============================================================================
\section{Task 2 --- the bare invariant space is $2$-dimensional (the sanity anchor)}
%==============================================================================

\subsection{The $(E_{11},u)$-forced $\so(3,1)$ on the frame}
Mirroring \texttt{cartan\_phaseA\_coframe.vald02} (exact over $\Q$): from the
calibration anchors $\dim\mathrm{Stab}_{E_6}(E_{11})=61$ and
$\dim\mathrm{Stab}_{V_0}=45=\dim\mathrm{Spin}(9,1)$, the residual structure group is
the nullspace over $\Q$ of $\{\mathrm{Stab}_{V_0}\text{ preserving the }\mathbb{C}_u$
$4$-space$\}$, of dimension $21=\so(3,1)[6]\oplus\so(6)[15]$. Restricting to the
$\mathbb{C}_u$ $4$-space $[1,2,3,10]$ and taking the $6$ independent generators
$L^{(k)}$ gives the Lorentz block: each kills $G_{\rm DET2,RAW}$ ($A^{\!\top}G+GA=0$),
and in the orthonormal frame each kills $\eta$. These $6$ forced generators span
\emph{exactly} the canonical Lorentz algebra $\so(\eta)=\{A:A^{\!\top}\eta+\eta A=0\}$
($\dim=6$, $\dim(\so(\eta)+\text{forced})=6$), re-confirming Phase~75.

\subsection{The invariance nullspace}
A quadratic-in-curvature $4$-form is $T_{abcd}\,R^{ab}\wedge R^{cd}$ with $T_{abcd}$
antisymmetric in $(a,b)$ and in $(c,d)$ and symmetric under $(ab)\leftrightarrow(cd)$
(the two $2$-forms commute). This is a $21$-dimensional parameter space ($6$ antisym
pairs, a symmetric $6\times6$). Imposing $\so(3,1)$-invariance,
\begin{equation}
  (\delta_g T)_{abcd}=-\bigl[g^e{}_a T_{ebcd}+g^e{}_b T_{aecd}
  +g^e{}_c T_{abed}+g^e{}_d T_{abce}\bigr]=0
  \quad\forall\,g\in\so(3,1),
\end{equation}
the exact nullspace over $\Q$ has
\begin{equation}\boxed{\ \dim_{\rm bare}=2\ }\qquad(\text{Euler}+\text{Pontryagin}).\end{equation}

\begin{remark}[the anchor is a hard gate]
If $\dim_{\rm bare}\neq2$ the $\so(3,1)$ setup would be buggy --- the baseline
\emph{must} be Euler $+$ Pontryagin (standard $4$d topology; Nieh--Yan is torsionful,
excluded for the Levi--Civita $\om$). The driver halts before any verdict in that
case. Here $\dim_{\rm bare}=2$, so the audit proceeds.
\end{remark}

\subsection{Identification and independence}
The two known invariants
\begin{equation}
  \text{Euler: } T^{\rm E}_{abcd}=\varepsilon_{abcd},\qquad
  \text{Pontryagin: } T^{\rm P}_{abcd}=\eta_{ac}\eta_{bd}-\eta_{ad}\eta_{bc},
\end{equation}
both lie in the bare invariant space, and $\operatorname{rank}\{T^{\rm E},T^{\rm
P}\}=2$ --- they are independent and \emph{span} the $2$-dimensional space.
Contracting them with two generic algebraic-curvature tensors (full Riemann
pair-symmetries) gives two-by-two value matrices of rank $2$ (e.g.\ on one generic
curvature, Euler $=-2528/49$, Pontryagin $=416/49$), certifying they are distinct,
linearly-independent invariants. (The Phase-77 \emph{diagonal} warped-reference
$R[\om]$ --- a genuine Lorentz-block Riemann of $g=\eu\cdot\eu$ but curved in a single
$2$-plane --- gives Pontryagin $\equiv0$, the exact and expected fact that the
Pontryagin density vanishes on diagonal metrics; the independence is therefore
certified on the generic curvature, not on this special $R$.)

\medskip\noindent\emph{Conclusion of Task~2.} Among quadratic-in-curvature
invariants $\varepsilon$ is \emph{not} unique: the bare space is already
$2$-dimensional. Forced-ness now hinges on whether the trace-form restriction
collapses $2\to1$ \emph{and} singles out $\varepsilon$ over Pontryagin \emph{and}
fixes the normalization.

%==============================================================================
\section{Task 3 --- the decisive trace-form-invariant subspace}
%==============================================================================
The decisive question is which of the bare invariants are buildable from the
\emph{intrinsic} data alone: the symmetric bilinear $\eta=\Tr|_{\rm frame}$
(Lemma~\ref{lem:eta}) and the symmetric trilinear $\det\nolimits_3$ (its polarization
$d$), with the $(E_{11},u)$-forced structures, and \emph{no} posited $\varepsilon$.

\subsection{Pontryagin is reachable; $\varepsilon$ is not a tensorial function of $\eta$}
The Pontryagin tensor $\eta_{ac}\eta_{bd}-\eta_{ad}\eta_{bc}$ is a pure $\eta$-product;
$\mathrm{span}\{T^{\rm P}\}$ has dimension $1$. The remaining $\eta$-product
$\eta_{ab}\eta_{cd}$ is symmetric in $(ab)$ and dies under the antisymmetrization, so
the $\eta$-tensorial invariant subspace is exactly $\mathrm{span}\{T^{\rm P}\}$. The
totally antisymmetric $\varepsilon_{abcd}$ is \emph{not} a polynomial in the symmetric
$\eta$ (the driver confirms $\varepsilon\notin\mathrm{span}\{T^{\rm P}\}$): no
contraction of symmetric tensors yields the antisymmetric orientation tensor.

\subsection{The cubic norm supplies no orientation: $\det_3$ vanishes on the Lorentz block}
\begin{proposition}[$\det\nolimits_3|_{\mathbb{C}_u\text{-block}}\equiv0$]\label{prop:det3}
The restriction of $\det\nolimits_3$ (equivalently its polarization $d$) to the four
soldered $V_0\cong\R^{3,1}$ frame directions $\{\beta,\gamma,p,q\}=[1,2,3,10]$ is
identically zero. Concretely, for the $\hh_3(\OO)$ layout the cubic norm is
\begin{equation}
  N(X)=\alpha\beta\gamma-\alpha|x_1|^2-\beta|x_2|^2-\gamma|x_3|^2+2\,\mathrm{Re}((x_2x_1)x_3),
\end{equation}
and on the block $\{\beta,\gamma,p=\mathrm{Re}(x_1),q=\langle x_1,e_7\rangle\}$ every
surviving term needs $\alpha$ (the $|x_1|^2$ term multiplies $\alpha$, not $\beta$;
the $\alpha\beta\gamma$ term needs $\alpha$), and $\alpha$ is \emph{outside} the
block. Hence $N|_{\rm block}\equiv0$, and the symmetric trilinear $d|_{\rm block}=0$
in all $64$ components (exact over $\Q$).
\end{proposition}
A symmetric trilinear that \emph{vanishes} on the block supplies no symmetric metric,
no antisymmetric $2$-form, no Pfaffian, and no orientation there. The $\det\nolimits_3$
route to $\varepsilon$ (route (b): a polarization/orientation on the Peirce block) is
therefore closed. The only orientation available is the metric volume form.

\subsection{$\varepsilon$ via the metric volume form (route a): reachable but non-unique and orientation-ambiguous}
A non-degenerate metric fixes a volume form up to orientation:
$\varepsilon_{abcd}=\sqrt{|\det g|}\,[abcd]$. In the orthonormal frame
$\det\eta=-1$, so $\sqrt{|\det\eta|}=1$: the $\varepsilon$ tensor is numerically
reachable from $\eta$ \emph{provided} an orientation (a discrete sign $[abcd]$) is
chosen. The normalization $\sqrt{|\det\eta|}=1$ is fixed by $\eta$; the orientation
($\varepsilon$ vs $-\varepsilon$) is \emph{not} fixed by the symmetric $\eta$ alone
--- it is precisely the ``broken by hand'' choice. Admitting the metric volume form,
$\mathrm{span}\{T^{\rm P},\varepsilon\}$ has dimension $2$: \emph{both} Pontryagin and
$\varepsilon$ are reachable, so $\varepsilon$ is one choice among several.

\subsection{The decisive dimension and the normalization}
Collecting:
\begin{align}
  \dim(\text{trace-form-invariant subspace}) &= 1
  &&(\eta\text{-tensorial only: }\mathrm{span}\{T^{\rm P}\};\ \varepsilon\text{ unreachable}),\\
  &= 2 &&(\text{admitting the volume form: }\mathrm{span}\{T^{\rm P},\varepsilon\};\ \varepsilon\text{ non-unique}).
\end{align}
Neither reading is a STRONG WIN. A strictly $1$-dimensional subspace whose generator
is $\varepsilon$ (not Pontryagin) with a $\det\nolimits_3$-fixed normalization would
be required; instead the $1$-dimensional $\eta$-tensorial subspace is Pontryagin
(\emph{not} $\varepsilon$), and admitting $\varepsilon$ makes the subspace
$2$-dimensional. As for the gravitational normalization ($\to 1/16\pi G$): by
Proposition~\ref{prop:det3} $\det\nolimits_3$ vanishes on the Lorentz block, so it
fixes \emph{nothing} about the $\varepsilon$-contraction's scale; the Einstein
coupling is free/imported.

\subsection{The decisive triple (input to the 78-02 verdict ladder)}
\begin{equation}\boxed{
\begin{array}{l}
\dim(\text{trace-form-invariant subspace}) = 1\ (\eta\text{-tensorial})\ /\ 2\ (\text{with volume form}),\\[2pt]
\varepsilon\text{-in-span} = \text{YES, only via the metric volume form (orientation a discrete,}\\
\qquad\text{non-intrinsic choice; }\det\nolimits_3\text{ supplies no }\varepsilon\text{); not singled out over Pontryagin},\\[2pt]
\text{normalization} = \text{free/imported }(\det\nolimits_3\equiv0\text{ on the Lorentz block}).
\end{array}}
\end{equation}
In every reading the STRONG-WIN condition (\,$\dim=1$ \emph{and} generator
$=\varepsilon$ \emph{and} normalization $\det\nolimits_3$-fixed, all three\,) fails.
The decisive input to Plan~78-02 is therefore $\texttt{fp-imported-action}$ --- the
high / most-likely outcome --- stated here at true strength: the MM
$\varepsilon$-contraction (and its normalization) is \emph{not} forced by the
$\hh_3(\OO)$ trace form / cubic norm; the symmetric intrinsic data give the symmetric
Pontryagin invariant for free but never single out the antisymmetric $\varepsilon$
with a fixed coupling. The Einstein term, recoverable only from a posited
$\int\varepsilon F\wedge F$, is an imported action.

%==============================================================================
\section{Task 4 (Plan 78-02) --- the verdict ladder applied to the triple}
%==============================================================================
The driver implements a \emph{deterministic} map $\texttt{verdict}(\text{triple})$ that
reads the exact-over-$\Q$ triple of \S3.6 and returns the categorical milestone
verdict via the ladder below, with no new computation, no posited action, and no
float. The map is \emph{not} hard-wired to a negative: on a hypothetical genuinely
forced triple (a strictly $1$-dimensional trace-form-invariant subspace whose
\emph{generator is} $\varepsilon$, with an intrinsic $\det\nolimits_3$-supplied
orientation and a $\det\nolimits_3$-fixed normalization) it returns STRONG WIN, and
breaking any \emph{single} clause flips the output to $\texttt{fp-imported-action}$ ---
verified in the driver's unit tests, so the verdict is genuinely \emph{read off} the
triple, not asserted.

\begin{definition}[the verdict ladder]\label{def:ladder}
Declare \textbf{STRONG WIN} iff all three clauses hold, each exact over $\Q$:
\begin{enumerate}
\item[(C1)] the trace-form-invariant subspace is strictly $1$-dimensional \emph{with
its single generator equal to} the $\varepsilon$-contraction (i.e.\ $\varepsilon$ is
singled out over Pontryagin by the intrinsic data);
\item[(C2)] $\varepsilon$ is a \emph{genuine intrinsic} invariant --- its orientation
fixed by $\Tr/\det\nolimits_3$, not a discrete by-hand choice (the metric-volume-form
route, whose orientation sign is put in by hand, does \emph{not} satisfy this);
\item[(C3)] the $\varepsilon$-contraction's gravitational normalization is fixed by the
cubic norm $\det\nolimits_3$.
\end{enumerate}
Any shortfall $\Rightarrow$ \textbf{\texttt{fp-imported-action}}, with the shortfall
named explicitly.
\end{definition}

\begin{proposition}[the milestone verdict, at true strength]\label{prop:verdict}
Applying Definition~\ref{def:ladder} to the decisive triple of \S3.6, the verdict is
$\boxed{\texttt{fp-imported-action}}$, on all three counts:
\begin{itemize}
\item \textbf{(C1) fails}: the $\eta$-tensorial subspace is $1$-dimensional but its
generator is \emph{Pontryagin}, not $\varepsilon$; admitting the metric volume form
makes the subspace $2$-dimensional ($\mathrm{span}\{T^{\rm P},\varepsilon\}$), so
$\varepsilon$ is one choice among several. There is no reading in which a strictly
$1$-dimensional subspace is generated by $\varepsilon$.
\item \textbf{(C2) fails}: $\varepsilon$ is reachable \emph{only} as
$\sqrt{|\det\eta|}\,\varepsilon$, whose orientation ($\varepsilon$ vs $-\varepsilon$) is
a discrete choice not fixed by the symmetric $\Tr/\det\nolimits_3$ data; by
Proposition~\ref{prop:det3} $\det\nolimits_3$ supplies no $\varepsilon$.
\item \textbf{(C3) fails}: the Einstein normalization ($\to 1/16\pi G$) is free/imported
--- $\det\nolimits_3$ vanishes identically on the soldered Lorentz block, so it pins no
scale; the coupling comes with the posited action.
\end{itemize}
The symmetric intrinsic data ($\Tr$, $\det\nolimits_3$) give the symmetric Pontryagin
invariant for free but never single out the antisymmetric $\varepsilon$ with a fixed
coupling. The Einstein term, recoverable only from a posited
$\int\varepsilon F\wedge F$, is an imported action.
\end{proposition}

\paragraph{This is the high / most-likely outcome --- reported at true strength.} The
contract registers \texttt{fp-imported-action} as the high-probability verdict (the
bare invariant space is already $2$-dimensional, and the symmetric-data-cannot-build-
antisymmetric-$\varepsilon$ obstruction is structural). The negative-result-is-success
discipline cuts \emph{both} ways: we neither inflate this honest partial into a win, nor
deflate a genuine win out of conservatism. Here the exact-over-$\Q$ triple genuinely
fails all three STRONG-WIN clauses, so \texttt{fp-imported-action} \emph{is} the verdict
at true strength. (Had the triple shown $\dim=1$ with $\varepsilon$ as the generator and
a $\det\nolimits_3$-fixed normalization, Proposition~\ref{prop:verdict} would have read
STRONG WIN --- the map is symmetric, as the driver's forced-triple unit test confirms.)

%==============================================================================
\section{The $\Lambda=0$ corollary --- a second, independent argument}
%==============================================================================
A second, logically independent argument reinforces \texttt{fp-imported-action} and is
read directly off the $\varepsilon F\wedge F$ expansion~\eqref{eq:ehexpand} together
with the Phase-77 measurement.

\begin{proposition}[$\Lambda=0$ degeneration to Gauss--Bonnet]\label{prop:lambda0}
In~\eqref{eq:ehexpand} the Einstein--Hilbert term carries coefficient
$-\tfrac{2\Lambda}{3}$ (linear in $\Lambda$) and the cosmological term carries
$\tfrac{\Lambda^2}{9}$ (quadratic in $\Lambda$), while the Gauss--Bonnet term is
$\Lambda$-independent. Phase~77 \emph{measured} $\Lambda=0$ at the $M=0$ vacuum
(CONVENTIONS~\S6); hence at the measured vacuum both the EH and the cosmological terms
\emph{vanish}, and
\begin{equation}
  \varepsilon_{abcd}F^{ab}\wedge F^{cd}\ \xrightarrow{\ \Lambda=0\ }\
  \varepsilon_{abcd}R^{ab}\wedge R^{cd}\quad(\text{pure Gauss--Bonnet / Euler}),
\end{equation}
a \emph{topological} invariant (the Euler density) yielding no equations of motion and
no general relativity. So even if the $\varepsilon$-action were posited, at the measured
vacuum it is topological, not gravitational. We do \emph{not} reintroduce $\Lambda<0$ /
$\R\times H^3$ (that thesis was falsified in v17.0 / Phase~70.1).
\end{proposition}

This is a structural statement read off the $\Lambda$-formal expansion plus the
Phase-77 measurement; it introduces no new approximation and no posited action. It is
independent of the invariant-count argument: the count shows $\varepsilon$ is not forced
\emph{as an invariant tensor}; the $\Lambda=0$ corollary shows that even a posited
$\varepsilon$-action would not give gravity \emph{at the measured vacuum}.

%==============================================================================
\section{The imported-action contrast class (GST / Singh / Castro)}
%==============================================================================
The verdict places this route in an explicit, well-populated equivalence class:
exceptional-algebra / unification gravity programs that recover Einstein gravity by
\emph{positing} an action and reading off the gravitational term from the matching
geometry. Stating precisely \emph{where} each imports its action makes ``posited by
hand'' concrete and shows why \texttt{fp-imported-action} (if that is the outcome ---
and here it is) is an honest partial in the same class, not a derivation.

\begin{itemize}
\item \textbf{GST (Gunaydin--Sierra--Townsend, 1983--84) magic $\mathcal{N}=2$ MESGT on
$E_{6(-26)}/F_4$.} \emph{Imports:} the $\mathcal{N}=2$ Maxwell--Einstein supergravity
\emph{Lagrangian} is posited; the $-R/2$ Einstein term is the assumed $\mathcal{N}=2$
multiplet's own output --- the cubic $C_{IJK}$ that defines the special-real geometry
\emph{also} fixes the gravitational term, so Einstein gravity is built in with the
supermultiplet, not derived from the algebra's invariant geometry. This is the
in-program $\texttt{fp-imported-action}$ precedent (the dead $47\text{--}50,\,53$
route). We cite GST for the $E_{6(-26)}/F_4$ geometry/contrast \emph{only}; we do
\emph{not} adopt its Lagrangian --- that posited Lagrangian is exactly the
\texttt{fp} under audit.
\item \textbf{Singh \emph{et al.} (2020--, arXiv:2009.05574) trace dynamics / division
algebras.} \emph{Imports:} a Planck-scale trace-dynamics Lagrangian together with the
Connes spectral action principle; gravity is put in via the spectral/``aikyon''
action, not derived from the algebra's invariant geometry. A contemporary
exceptional-Jordan-gravity program in the same imported-action class.
\item \textbf{Castro Perelman (2000s--) $E_8$ / Clifford gravity.} \emph{Imports:} a
candidate gravitational action for an $E_8$ gauge theory of gravity in $D=8$ Clifford
($Cl(16)$) space (conformal-gravity / Yang--Mills GUT shape); the gravitational action
is constructed/posited, not forced by a cubic-norm invariant count.
\end{itemize}

\paragraph{Our forced-vs-posited TEST, positioned against the class.} All three
programs posit an action and read Einstein off the matching geometry. Our audit is
non-circular precisely because it does \emph{not} posit $\int\varepsilon F\wedge F$: it
asks whether the $\varepsilon$-contraction is forced by $\Tr(X\circ Y)$ and
$\det\nolimits_3$. The milestone's hopeful hypothesis was that $\hh_3(\OO)$'s
\emph{specific} cubic-norm / trace-form structure --- richer than a generic
$\mathrm{SO}(3,1)$ frame --- might single out $\varepsilon$ where generic MM does not
(the exceptional-structure escape hatch). The exact-over-$\Q$ count
\emph{closes that hatch}: $\det\nolimits_3$ vanishes on the Lorentz block and the
symmetric $\Tr$ cannot produce the antisymmetric $\varepsilon$, so $\hh_3(\OO)$ does
\emph{not} force $\varepsilon$. By Proposition~\ref{prop:verdict} this route is therefore
in the \emph{same honest-partial class} as GST/Singh/Castro: curved, matter-sourced, and
carrying a forced $\mathrm{SO}(3,1)$ coframe (Phases~75/77), but with Einstein structure
obtainable only via a posited action. It is a publishable closure --- a clean negative
that says precisely what ``posited'' means --- not a derivation of gravity.

%==============================================================================
\section{Milestone-closing summary (v18.0) and the final verdict}
%==============================================================================

\subsection{The v18.0 through-line (Phases 74--78)}
v18.0 asked whether gravity is the curvature of the Peirce-frame Cartan / MacDowell--
Mansouri connection on $\hh_3(\OO)$ --- the imaginary part of the self-model quantum
geometric tensor (the Berry / Lie-sector curvature), a \emph{different} tensor from the
v17.0 cone-Hessian (the real part), so the v17.0 NONE verdict does not bind it. The
phases, each a cheap kill/audit gate:
\begin{description}
\item[Phase 74 (Phase 0 --- foundation).] Certified the exact-over-$\Q$ foundation: the
det SSOT ($324/324=\dim F_4=52$), the soldering form $dE$ as $V_{1/2}$-valued, the
anchors $24/28/3$, $K=-\tfrac12$. \emph{Certified.}
\item[Phase 75 (Phase A --- coframe-reduction kill gate). SURVIVES.] $(E_{11},u=e_7)$
forces a $4$d Lorentzian $(1,3)$ coframe carrying $\mathrm{SO}(3,1)$:
$\dim\pi_u(V_{1/2})=4$; the soldering bilinear $B$ into $V_0\cong\R^{3,1}$ has signature
$(1,3)$ with a forward null cone; the residual $21=\so(3,1)[6]\oplus\so(6)[15]$ has the
$\so(6)$ a genuine ideal acting trivially on spacetime ($\mathrm{res}/\so(6)=\so(3,1)$),
so $\so(3,1)$ is \emph{forced}, not arbitrarily chosen.
\item[Phase 76 (Phase A.5 --- Berry curvature). Reinterpreted.] $F_B=\mathrm{Im(QGT)}$ is
the internal $\mathrm{SU}(4)/\so(6)$ gauge curvature, \emph{not} gravity; the proposed
soft kill tested the wrong object (a quadratic-in-$F$ Maxwell/Pontryagin stress, whose
tracelessness in $4$d is a tautology that also kills real GR). Gravity is the soldering
Riemann $R[\om]$ (the \emph{linear}-in-Riemann object). The conjunctive $75\wedge76$ gate
was retired; Phase~75 alone greenlit Phase~B.
\item[Phase 77 (Phase B --- full Cartan curvature). NEGATIVE / fp-imported-action
partial.] $F=dA+A\wedge A$ has Lorentz block $==$ the genuine Levi-Civita Riemann of
$g=\eu\cdot\eu$ ($6/6$ components exact over $\Q$, torsion $=0$); the $M=0$ vacuum is
flat ($\Lambda=0$ measured); but $G[g]$ is \emph{not} Einstein-form $\kappa T+\Lambda g$
for any single global $(\kappa,\Lambda)$ against an AST-guarded order-matched independent
$T[M]$ (a two-axis failure: per-point $\Lambda$ varies and the tensor support
mismatches). Curved and matter-sourced, but not Einstein-structured \emph{without} a
posited action. Independent of the v17.0 Ph73 NONE (different tensor).
\item[Phase 78 (Phase C --- circularity audit; this work). \texttt{fp-imported-action}.]
The MM $\varepsilon$-contraction (and its normalization) is \emph{not} forced by the
$\hh_3(\OO)$ trace form / cubic norm (Proposition~\ref{prop:verdict}); the
$\Lambda=0$ corollary (Proposition~\ref{prop:lambda0}) reinforces this.
\end{description}

\subsection{The combined milestone verdict}
The two decisive phases settle the two faces of the same question and agree:
\begin{itemize}
\item Phase~77 settled the \emph{dynamical} side: the assembled curvature is genuine
Riemann but the field equations $G[g]$ are not intrinsically Einstein --- Einstein
structure, if any, requires a posited action.
\item Phase~78 settled the \emph{tensor-level / action} side: the $\varepsilon$-action
that \emph{would} supply Einstein is itself not forced by the intrinsic $\Tr/\det_3$
data --- it is imported.
\end{itemize}
Together: \textbf{v18.0 closes \texttt{fp-imported-action}}. The $\hh_3(\OO)$
Peirce-frame Cartan / MM connection yields a curved, matter-sourced, position-dependent
Lorentzian slice carrying a forced $\mathrm{SO}(3,1)$ coframe, but \emph{not} Einstein
gravity without a posited action --- an honest partial in the same equivalence class as
GST / Singh / Castro, and a full publishable closure of the milestone.

\paragraph{Scope (stated honestly).} The invariant-count verdict of
Proposition~\ref{prop:verdict} certifies whether $\varepsilon$ is forced \emph{as an
invariant tensor}; by itself it does not prove the \emph{dynamical} claim ``gravity is
the Lie-sector connection curvature.'' That dynamical claim was settled NEGATIVE in
Phase~77 ($G[g]$ not Einstein-form intrinsically). The two phases \emph{together} give
the full milestone verdict; neither alone does. The chief modeling choice is that the
``forced'' criterion is operationalized at the \emph{broken} $\so(3,1)$ level --- the
$\mathrm{Spin}(9,1)\to\mathrm{SO}(3,1)$ breaking is exactly the audited ``by hand''
step; a reviewer counting at the full $\mathrm{SO}(4,1)/\mathrm{SO}(3,2)$ level could
differ, but the contract, ROADMAP, and the ``by hand'' locus all specify the broken
$\so(3,1)$ count.

\paragraph{Disconfirming observations (what would have changed the verdict).} The
verdict is falsifiable and rests on exact-over-$\Q$ facts: (i) had $\dim_{\rm bare}\neq
2$ the $\so(3,1)$ setup would be buggy --- it is $2$ (triple-confirmed), so the baseline
is sound; (ii) had $\det\nolimits_3$ supplied a genuine orientation / Pfaffian on the
Lorentz block, $\varepsilon$ could have been singled out --- it does not
($\det\nolimits_3\equiv0$ there); (iii) had the $\varepsilon$ normalization been read off
$\det\nolimits_3$, the $\dim=1$ reading could approach a win --- it is free/imported; and
(iv) had ``$\varepsilon$-contraction gives Einstein--Hilbert'' been offered as the forced
witness, that would be the posited action smuggled in --- which we counted, did not
posit. None of (ii)--(iv) materialized, so the verdict is \texttt{fp-imported-action} at
true strength.

\medskip\noindent\textbf{Final verdict (Plan 78-02, ratified at the blocking
checkpoint):} $\boxed{\texttt{fp-imported-action}}$ --- the MM $\varepsilon$-contraction,
hence the Einstein--Hilbert term, is not forced by the $\hh_3(\OO)$ trace form / cubic
norm; it is recoverable only from a posited action, and at the Phase-77-measured
$\Lambda=0$ vacuum even that posited action is topological. This closes
\texttt{claim-forced-einstein} and milestone v18.0.

\end{document}
